\section{Public Key Encryption}

One encryption scheme that can be derived from
DH is El Gamal.
The El Gamal Encryption Scheme is as follows:
\begin{itemize}
    \item \texttt{Gen}: on Group $\mathbb{G}$ of order $q$
          and generator $g$, choose a uniform $x \in \mathbb{Z}_q$
          and compute $h = g^x$. Public key is $(\mathbb{G}, q, g, h)$
          and the private key is $(\mathbb{G}, q, g, x)$.
    \item \texttt{Enc}: on input public key $pk = (\mathbb{G}, q, g, h)$
          and message $m \in \mathbb{G}$ choose uniform $y \in \mathbb{Z}_q$
          and output ciphertext $g^y, h^y\cdot m$.
    \item \texttt{Dec}: on input private key $sk = (\mathbb{G}, q, g, x)$
          and ciphertext $(c_1, c_2)$ output $m = \frac{c_2}{c_1^x}$.
\end{itemize}
El Gamal encryption is, nicely, multiplicatively homomorphic.

An encryption scheme derived from RSA follows:
\begin{itemize}
    \item \texttt{Gen}: run \texttt{GenRSA} to get $(N, e, d)$,
          public key is $⟨N, e⟩$ and private key is $⟨N, d⟩$
    \item \texttt{Enc}: on input public key $pk = ⟨N, e⟩$ and message
          $m \in \mathbb{Z}_N^*$, compute ciphertext $c = [m^e \mod N]$
    \item \texttt{Dec}: on input private key $sk = ⟨N, d⟩$ and
          ciphertext $c \in \mathbb{Z}_N^*$ compute the
          message $m = [c^d \mod N ]$
\end{itemize}

